% !TEX root =../ECE343_Slides_Rogers.tex


%%%%%%%%%%%%%%%%%%%%%%%%%%%%%%%%%%%%%%%%%%%%%%%%%%%%%%
%%%%%%%%%%%%%%%%%%%%%%%%%%%%%%%%%%%%%%%%%%%%%%%%%%%%%%

%\begin{comment}
\section{Schedule}
\begin{frame}{Schedule}
	\def\myfontsize{footnotesize}
	\setcounter{mylnum}{24}
%\def\myfontsize{normalsize}
\begin{center}
\begin{\myfontsize}
\begin{tabular}{| c | l | c | l |}
\hline
\textbf{Lesson} 		  	& \textbf{Topics}             & \pbox{5cm}{\textbf{Reading}}	& \textbf{Notes}		\\ \hline\hline
\ifnum\lessonNum=\value{mylnum} \rowcolor{green!20!white}\fi
	\themylnum  
	& BVP (Capacitor, Coax), Electric Potential	
	& 5.16
	&  LO 5.3  
\stepcounter{mylnum}\\ \hline \hline
\ifnum\lessonNum=\value{mylnum} \rowcolor{green!20!white}\fi
	\themylnum 	
	& Electric Boundary Conditions
	& 5.18-5.22	
	& LO 5.4			
\stepcounter{mylnum}\\ \hline
\ifnum\lessonNum=\value{mylnum} \rowcolor{green!20!white}\fi
	\themylnum 	
	& BC's/Capacitance			
	& 5.23-5.24
	& 	
\stepcounter{mylnum}\\ \hline
\ifnum\lessonNum=\value{mylnum} \rowcolor{green!20!white}\fi
	\themylnum 	
	& Electrostatic Lab			
	& Lab packet
	& 	
\stepcounter{mylnum}\\ \hline
\rowcolor{gray!20!white}\multicolumn{4}{| c |}{\textbf{Chapter 6: Steady Current \& Conductivity}}	\\ \hline\hline
\ifnum\lessonNum=\value{mylnum} \rowcolor{green!20!white}\fi
	\themylnum 	
	& Current density, Conductivity, Resistance
	& Ch. 6
	&   					
\stepcounter{mylnum}\\ \hline\hline
\rowcolor{gray!20!white}\multicolumn{4}{| c |}{\textbf{Chapter 7: Magnetostatics}}						\\ \hline\hline
\ifnum\lessonNum=\value{mylnum} \rowcolor{green!20!white}\fi
	\themylnum 	
	&  Magnetic Fields, Gauss' Law for Magnetism	
	& 7.1-7.3		
	&LO 7.1-7.2 \textbf{GR 2} 				
\stepcounter{mylnum}\\ \hline
\ifnum\lessonNum=\value{mylnum} \rowcolor{green!20!white}\fi
	\themylnum 	
	& Ampere's Law 	
	& 7.4-7.5, 7.9
	& LO 7.3		
\stepcounter{mylnum}\\ \hline 
\ifnum\lessonNum=\value{mylnum} \rowcolor{green!20!white}\fi
	\themylnum 	
	&  Ampere's Law Applied
	&  7.6-7.8
	& LO 7.3
\stepcounter{mylnum}\\ \hline 
\ifnum\lessonNum=\value{mylnum} \rowcolor{green!20!white}\fi
	\themylnum	
	& Magnetic Boundary Conditions	
	& 7.10-7.14	
	& LO 7.4 		
\stepcounter{mylnum}\\ \hline 
\ifnum\lessonNum=\value{mylnum} \rowcolor{green!20!white}\fi
	\themylnum  
	& Magnetostatics Lab
	& Lab packet
	&  		
\stepcounter{mylnum}\\ \hline \hline
\rowcolor{gray!20!white}\multicolumn{4}{| c |}{\textbf{Chapter 8: Time-Varying Fields}}					\\ \hline\hline
\ifnum\lessonNum=\value{mylnum} \rowcolor{green!20!white}\fi
	\themylnum 	
	& Faraday's Law, Displacement Current, Maxwell Eqns
	& 8.3, 8.8, 8.9
	& LO 8.1
\stepcounter{mylnum}\\ \hline 
\ifnum\lessonNum=\value{mylnum} \rowcolor{green!20!white}\fi
	\themylnum	
	& Faraday's Law, Displacement Current, Maxwell Eqns 
	&  8.3, 8.8, 8.9
	& LO 8.2-8.3
\stepcounter{mylnum}\\ \hline 

\rowcolor{gray!20!white}\multicolumn{4}{| c |}{\textbf{Chapter 9: Plane Waves in Lossless Media}}					\\ \hline\hline
\ifnum\lessonNum=\value{mylnum} \rowcolor{green!20!white}\fi
	\themylnum	
	& Maxwell's Equations in lossless, source-free region		
	& 9.1-9.2
	& LO 9.1	 
\stepcounter{mylnum}\\ \hline
\ifnum\lessonNum=\value{mylnum} \rowcolor{green!20!white}\fi
	\themylnum	
	& Maxwell's Equations in lossless, source-free region		
	& 9.1-9.2
	& LO 9.1	 
\stepcounter{mylnum}\\ \hline
%\ifnum\lessonNum=\value{mylnum} \rowcolor{green!20!white}\fi
%	\themylnum	
%	& Flex day
%	& 
%	& \textbf{GR 3}			 					
%\stepcounter{mylnum}\\ \hline 
\ifnum\lessonNum=\value{mylnum} \rowcolor{green!20!white}\fi
	\themylnum	
	& Uniform Plane Waves \& Polarization
	& 9.4-9.6
	& LO 9.2-9.3			
\stepcounter{mylnum}\\ \hline 
\ifnum\lessonNum=\value{mylnum} \rowcolor{green!20!white}\fi
	\themylnum	
	& Normal Incidence
	& Vol 2, 5.6-5.7
	&  LO 9.4-9.5
\stepcounter{mylnum}\\ \hline 
\ifnum\lessonNum=\value{mylnum} \rowcolor{green!20!white}\fi
	\themylnum	
	& Oblique Incidence			
	& Vol 2, 5.6-5.7
	& LO 9.5
\stepcounter{mylnum}\\ \hline 
\ifnum\lessonNum=\value{mylnum} \rowcolor{green!20!white}\fi
	\themylnum	
	& Oblique Incidence			
	& Vol 2, 5.6-5.7
	& LO 9.5
\stepcounter{mylnum}\\ \hline 
\ifnum\lessonNum=\value{mylnum} \rowcolor{green!20!white}\fi
	\themylnum	
	& Incidence Lab			
	& TBD
	& LO 9.5
\stepcounter{mylnum}\\ \hline 
%\ifnum\lessonNum=\value{mylnum} \rowcolor{green!20!white}\fi
%	\themylnum	
%	& Wrapup
%	& 
%	& 		 					
%\stepcounter{mylnum}\\ \hline 

\end{tabular}			
\end{\myfontsize}
\end{center}

\end{frame}

%%%%%%%%%%%%%%%%%%%%%%%%%%%%%%%%%%%%%%%%%%%%%%%%%%%%%
%%%%%%%%%%%%%%%%%%%%%%%%%%%%%%%%%%%%%%%%%%%%%%%%%%%%%

\def\obja{ I understand the process of decoupling and finding solutions for Maxwell's equations for dynamic fields.}
\def\objb{ I can articulate what a Uniform Plane Wave (UPW) is and describe key properties that relate the electric and magnetic field. }
\def\objc{ I can calculate the electric and magnetic fields for a UPW.}
\def\objd{ I can articulate the concept of polarization and determine the polarization for various UPWs.}
\def\obje{ I can determine the reflection and transmission of UPW at a boundary for Transverse Electric (TE) and Transverse Magnetic (TM) waves, using boundary conditions.}
\begin{frame}{Ch 9 Learning Objectives}
	\begin{textblock}{14.5}(0.5,2.5)
	\begin{myblock}[0.95]{}
	\begin{small}
	\begin{enumerate}[{9}.1]
		\ifnum\objNum=1
			\item \alert{ \obja }
		\else 
			\item \obja
		\fi 
		\ifnum\objNum=23
			\item \alert{ \objb }
			\item \alert{ \objc }
		\else 
			\item \objb
			\item \objc
		\fi 
		\ifnum\objNum=4
			\item \alert{ \objd }
		\else 
			\item \objd
		\fi 
		\ifnum\objNum=5
			\item \alert{ \obje }
		\else 
			\item \obje
		\fi 
	\end{enumerate}
	\end{small}
	\end{myblock}
	\end{textblock}
\end{frame}

%%%%%%%%%%%%%%%%%%%%%%%%%%%%%%%%%%%%%%%%%%%%%%%%%%%%%%
%%%%%%%%%%%%%%%%%%%%%%%%%%%%%%%%%%%%%%%%%%%%%%%%%%%%%%
\section{Maxwell's Equations}
\begin{frame}{Maxwell's Equations}
\begin{textblock}{14.5}(0.5,2.25)

\begin{center}
\begin{tabular}{| l | l |}
\hline
\textbf{Integral Form} & \textbf{Point Form}  \\\hline
 ${\displaystyle \iint \epsilon\vec{E}\mycdot \vec{dS}} = Q_{\text{enc}} $ 
 & $\nabla\mycdot\vec{E}  = \dfrac{\rho_v}{\epsilon}$   \\ \hline
${\displaystyle\oint \vec{H}\mycdot\vec{dS} } = 0 $ 
& $\nabla\mycdot\vec{H} = 0$ \\ \hline
$\begin{aligned}
\displaystyle\oint_C \vec{E}\mycdot d\ell &=-\partd{\Psi_{\text{m}}}{t} \\
\displaystyle\oint_C \vec{E}\mycdot d\ell &=  -\iint\partd{\vec{B}}{t}\mycdot \vec{dS}
\end{aligned}$
&$\begin{aligned}
\nabla\times\vec{E} &= -\mu\partd{\vec{H}}{t} \\
%\nabla\times\myhat{E}_s &= -j\omega\mu\myhat{H}_s
\end{aligned}$
\\ \hline
$\begin{aligned}
\displaystyle\oint_C\vec{H}\mycdot d\ell &= I_{\text{encl}} \\ 
\displaystyle\oint_C\vec{H}\mycdot d\ell &= \int \vec{J}\mycdot\vec{dS} 
 + \epsilon\partd{}{t}\iint\vec{E}\mycdot \vec{dS}
\end{aligned} $
& $\begin{aligned}
\nabla\times\vec{H} &= \sigma \vec{E} + \epsilon\partd{\vec{E}}{t} \\ 
%\nabla\times\myhat{H}_s &= \lp \sigma + j\omega \varepsilon \rp \myhat{E}_s
 \end{aligned}$ \\ \hline
\end{tabular}
\end{center}

\end{textblock}
\end{frame}

%%%%%%%%%%%%%%%%%%%%%%%%%%%%%%%%%%%%%%%%%%%%%%%%%%%%%%
%%%%%%%%%%%%%%%%%%%%%%%%%%%%%%%%%%%%%%%%%%%%%%%%%%%%%%
\section{Phasors}
\begin{frame}{Maxwell's Equations: Phasors}
\begin{textblock}{8}(0.5,3)
\circled{1}\; Starting with Gauss' Law:
\begin{small}
\visible<0|handout:1>{
\eq{
		\mydiv{\vec{E}} = \dfrac{\rho_v}{\varepsilon}
} % end eq
} % end visible
\end{small}
\circled{2}\; Assume that the electric and magnetic field are sinusoidally varying:
\vspace{-8pt}
\visible<0|handout:1>{
\eq{
	\vec{E},\;\vec{H} = A_{e,h} \cosp{\omega t - \beta z + \phi}
} % end eq
} % end visible

\vspace{-8pt}

\circled{3}\; Using a phasor representation of electric flux density and charge density:

\vspace{-18pt}

\begin{small}
\visible<0|handout:1>{
\eq{
	\vec{E} &= \text{Re} \left\{ \widetilde{\boldmath{E}} e^{j\omega t}  \right\} \\
	\implies \widetilde{\boldmath{E}} &=  e^{-j\lp\beta z - \phi\rp} = e^{-j\beta z}e^{j\phi} \\
	\rho_v &= \text{Re} \lbr \widetilde{\rho}_v e^{j\omega t}  \rbr
} % end eq
} % end visible
\end{small}
\end{textblock}
\begin{textblock}{7}(8.5,3)
\begin{small}
\circled{4}\; What is the phasor form of Gauss' Law?
\visible<0|handout:1>{
\eq{
	&\nabla \mycdot \lb \re{\widetilde{\boldmath{E}} e^{j\omega t} } \rb = \dfrac{1}{\varepsilon} \re{\widetilde{\rho}_v e^{j\omega t}}  \\[8pt]
	&\re{\nabla\mycdot \lb \widetilde{\boldmath{E}} e^{j\omega t} \rb } = \dfrac{1}{\varepsilon}\re{\widetilde{\rho}_v e^{j\omega t}}  \\[8pt]
	&\re{\lb \div{\widetilde{\boldmath{E}} }  \rb e^{j\omega t} } = \dfrac{1}{\varepsilon}\re{\widetilde{\rho}_v e^{j\omega t}}  
} % end eq
} % end visible
\end{small}
\vspace{-18pt}
\begin{myblock}[0.9]{Gauss' Law}
	\visible<0|handout:1>{
	\vspace{-12pt}
	\eq{
		\div{\widetilde{\boldmath{E}}} = \dfrac{\widetilde{\rho}_v}{\varepsilon}
		} % end eq
		\vspace{-12pt}
		} % end visible 
\end{myblock}
\end{textblock}
\end{frame}

%%%%%%%%%%%%%%%%%%%%%%%%%%%%%%%%%%%%%%%%%%%%%%%%%%%%%%
%%%%%%%%%%%%%%%%%%%%%%%%%%%%%%%%%%%%%%%%%%%%%%%%%%%%%%
\begin{frame}{Faraday's Law}
\begin{textblock}{8}(0.5,3)
\circled{1}\; Using the same approach, let's look at Faraday's Law:
\begin{small}
\visible<0|handout:1>{
\eq{
	\curl{\vec{E}} = -\partd{\vec{B}}{t} 
} % end eq
} % end visible 
\end{small}
\circled{2}\; Assume that the electric and magnetic field are sinusoidally varying:
\visible<0|handout:1>{
\vspace{-8pt}
\eq{
	\vec{E},\;\vec{H} = A_{e,h}\cosp{\omega t - \beta z + \phi}
} % end eq
} % end visible 

\vspace{-8pt}

\circled{3}\; Using the phasor representations: 

\vspace{-18pt}

\begin{small}
\visible<0|handout:1>{
\eq{
	\vec{E} &= \re{\phasor{E} e^{j\omega t} }   \\
	\vec{H} &= \re{\phasor{H} e^{j\omega t} }   \\
} % end eq
} % end visible
\end{small}
\end{textblock}
\begin{textblock}{7}(8.5,3)
\begin{small}
\circled{4}\; What is the phasor form of Faraday's Law?
\visible<0|handout:1>{
\eq{
	&\curl{\re{\phasor{E}}} = \re{-\partd{}{t}\phasor{B} e^{j\omega t}} \\[8pt]
	& \re{\lb \curl{\phasor{E}}\rb e^{j\omega t} } = \re{-\phasor{B} \partd{}{t} e^{j\omega t}}\\[8pt]
	&\re{\lb \curl{\phasor{E}}\rb e^{j\omega t} } = \re{\lb - j\omega \phasor{B}\rb e^{j\omega t}}
} % end eq
} % end visible 
\end{small}
\vspace{-18pt}
\begin{myblock}[0.9]{Faraday's Law}
\visible<0|handout:1>{
	\vspace{-12pt}
	\eq{
		 	\curl{\phasor{E}} = -j\omega\phasor{B}
		} % end eq
		\vspace{-12pt}
		} % end visible
\end{myblock}
\end{textblock}
\end{frame}
%%%%%%%%%%%%%%%%%%%%%%%%%%%%%%%%%%%%%%%%%%%%%
%%%%%%%%%%%%%%%%%%%%%%%%%%%%%%%%%%%%%%%%%%%%%

\begin{frame}{Maxwell's Equations: Phasor Forms}
\begin{textblock}{7}(0.5,2.5)
\begin{myblock}[0.95]{General}
\begin{small}
\vspace{-18pt}
	\eq{
		&\div{\phasor{E}} = \dfrac{\widetilde{\rho}_v}{\varepsilon} \\[8pt]
		&\div{\phasor{B}} = 0 \\[8pt]
		&\curl{\phasor{H}} = \phasor{J} + j\omega \varepsilon \phasor{E} = \lp \sigma + j\omega \varepsilon \rp\phasor{E}\\[8pt]
		&\curl{\phasor{E}} =  -j\omega\mu \phasor{H}
		}
\end{small}
\vspace{-12pt}
\end{myblock}
\end{textblock}

\begin{textblock}{7}(8,2.5)
\begin{myblock}[0.9]{Charge-free Region}
\begin{small}
	\vspace{-18pt} 
	\eq{
		&\div{\phasor{E}} =\handoutgap{1in}{  \dfrac{\widetilde{\rho}_v}{\varepsilon} }\\[8pt]
		&\div{\phasor{B}} = \handoutgap{1in}{ 0} \\[8pt]
		&\curl{\phasor{H}} =\handoutgap{1in}{ \lp \sigma + j\omega \varepsilon \rp\phasor{E} }\\[8pt]
		&\curl{\phasor{E}} =  \handoutgap{1in}{ -j\omega\mu \phasor{H}}
		}
\end{small}
\vspace{-12pt}
\end{myblock}
\end{textblock}

\begin{textblock}{14.5}(0.5,11.5)

	\eq{
		\curl{\phasor{H}} = \lp \underbrace{\sigma}_{\text{Real}} + \underbrace{j\omega\varepsilon}_{\text{imag}}\rp \phasor{E}  \implies  \underbrace{\tanp{\delta} = \dfrac{\sigma}{\omega \varepsilon} }_{\text{ Loss tangent}}
	}


\end{textblock}
\end{frame}

%%%%%%%%%%%%%%%%%%%%%%%%%%%%%%%%%%%%%%%%%%%%%
%%%%%%%%%%%%%%%%%%%%%%%%%%%%%%%%%%%%%%%%%%%%%

\section{Helmholtz Wave Equations}
\begin{frame}{Wave Equations}
\begin{textblock}{14.5}(0.5,3)
The phasor form of Maxwell-Faraday and Maxwell-Ampere's equations are \textbf{coupled}, \textbf{partial differential equations}. To find a solution:
\begin{itemize}
	\item Use phasors to find time-independent expressions
	\item Decouple equations 
	\item Solve resulting 2nd order ODEs using normal methods
	\item Interpret results
\end{itemize}

\begin{myblock}[0.9]{Phasor Forms}
\eq{
	\curl{\phasor{E}} = -j\omega \mu \phasor{H} \qquad
	\curl{\phasor{H}} = \lp \sigma + j\omega \varepsilon \rp \phasor{E}
}
\end{myblock}
\end{textblock}

\end{frame}

%%%%%%%%%%%%%%%%%%%%%%%%%%%%%%%%%%%%%%%%%%%%%
%%%%%%%%%%%%%%%%%%%%%%%%%%%%%%%%%%%%%%%%%%%%%

\begin{frame}{Decoupling Maxwell's Equations}
\begin{textblock}{7.5}(0.5,3.5)
Starting with Faraday's law:
\begin{small}
\visible<0|handout:1>{
\eq{
	&\curl{\phasor{E}} = -j\omega\mu \phasor{H} \\
	&\curl{\curl{\phasor{E}}} = \curl{\lp -j\omega\mu \phasor{H}\rp} \\
	&\nabla \lp \cancelto{0}{\div{\phasor{E}}}\rp - \mydyad{\nabla}^2\phasor{E} = -j\omega\mu\lp\curl{\phasor{H}}\rp \\
	&\text{Ampere:  } \curl{\phasor{H}} =  j\omega \varepsilon\phasor{E} \\[8pt]
	&\mydyad{\nabla}^2 \phasor{E} - \gamma^2\phasor{E} = 0\\
	&\gamma^2 =  j\omega\mu\lp\sigma + j\omega\varepsilon \rp 
} % end eq
} % end visible
\end{small}
\end{textblock}

\begin{textblock}{7}(8.5,3.5)
Your turn -- decouple starting with Ampere's Law...


\end{textblock}

\end{frame}

%%%%%%%%%%%%%%%%%%%%%%%%%%%%%%%%%%%%%%%%%%%%%
%%%%%%%%%%%%%%%%%%%%%%%%%%%%%%%%%%%%%%%%%%%%%

\begin{frame}{Helmholtz Wave Equations}
\begin{textblock}{7}(8.5,3.5)
\begin{myblock}[0.9]{}
	\eq{
		\handoutgap{0pt}{\mydyad{\nabla}^2 \phasor{E} - \gamma^2\phasor{E} = 0} \\[8pt]
		\handoutgap{0pt}{\mydyad{\nabla}^2 \phasor{H} - \gamma^2\phasor{H} = 0} \\[8pt]
		\handoutgap{0pt}{\gamma^2 =  j\omega\mu\lp\sigma + j\omega\varepsilon \rp }\\[8pt]
	}	
\end{myblock}
\end{textblock}
\end{frame}

